% =============================================================================
% NOTA EDITORIAL — EDITORIAL NOTE — 附录——编辑说明 (trilíngue)
% =============================================================================
\cleardoublepage
\thispagestyle{empty}
\chapter*{Apêndice --- Nota Editorial · Editorial Note · 附录——编辑说明}
\addcontentsline{toc}{chapter}{Apêndice · Editorial Note · 附录——编辑说明}
\markboth{Nota Editorial · Editorial Note · 编辑说明}{}

\begingroup

\noindent\textbf{Português.} Esta edição foi preparada a partir da transcrição pericial do caderno original do Paciente 1.260, encontrado em anexo ao prontuário do Hospital Colônia de Barbacena sob custódia do Instituto Médico Legal em 1979. O critério editorial adotado foi o de máxima fidelidade ao documento original, preservando sua estrutura fragmentária, suas lacunas, suas repetições e suas inconsistências internas. Os 78 fragmentos aqui reunidos seguem a ordem aproximada de escrita, reconstituída por datação relativa do suporte, da tinta e da caligrafia, conforme laudo pericial anexo (DOC-03).

\vspace{0.4em}
\noindent\textbf{Aparato paratextual:} \textbf{[N.E.]} --- Nota do Editor: comentário editorial sobre o conteúdo, a fonte ou os limites da reconstituição. \textbf{[ilegível]} --- trecho do original não decifrado. \textbf{[lacuna]} --- página ou trecho ausente. \textbf{[sic]} --- transcrição literal de erro ou idiossincrasia do original.

\vspace{0.4em}
\noindent\textbf{Limites:} esta edição não reivindica completude. O diário original continha aproximadamente 847 páginas; o que chega ao leitor é uma seleção editada, organizada em cinco dimensões temáticas (MEM, DOC, LUC, CONT) mais o Epílogo. A presente edição trilíngue apresenta cada fragmento em três línguas (Português → English → 简体中文), como espelho das edições publicadas.

\vspace{1em}
\noindent\rule{\textwidth}{0.5pt}
\vspace{0.8em}

\noindent\textbf{English.} This edition was prepared from the forensic transcription of the original notebook of Patient 1,260, found attached to the medical record of the Barbacena Hospital-Colony under the custody of the Legal Medical Institute in 1979. The editorial criterion adopted was maximum fidelity to the original document, preserving its fragmentary structure, its gaps, its repetitions, and its internal inconsistencies. The 78 fragments gathered here follow the approximate order of writing, reconstructed by relative dating of the support, the ink, and the handwriting, according to the attached forensic report (DOC-03).

\vspace{0.4em}
\noindent\textbf{Paratextual apparatus:} \textbf{[Ed. Note]} --- Editor's Note: editorial comment on the content, source, or limits of the reconstruction. \textbf{[illegible]} --- undeciphered passage. \textbf{[gap]} --- missing page or passage. \textbf{[sic]} --- literal transcription of an error or idiosyncrasy of the original.

\vspace{0.4em}
\noindent\textbf{Limits:} this edition does not claim completeness. The original diary contained approximately 847 pages; what reaches the reader is an edited selection, organized in five thematic dimensions (MEM, DOC, LUC, CONT) plus the Epilogue. This trilingual edition presents each fragment in three languages (Portuguese → English → Simplified Chinese), as a mirror of the published editions.

\vspace{1em}
\noindent\rule{\textwidth}{0.5pt}
\vspace{0.8em}

\noindent\textbf{中文.} 本版依据患者1,260原始笔记本的法证誊本编订，该笔记本于1979年由验尸机构保管，附于巴巴塞纳收容院病历。采用的编辑标准是对原始文件的最大忠实，保留其碎片化结构、空缺、重复和内部矛盾。这里汇集的78个碎片遵循大致的写作顺序，通过载体、墨迹和笔迹的相对年代测定重建，依据所附法证报告（DOC-03）。

\vspace{0.4em}
\noindent\textbf{副文本装置:} \textbf{[编者按]} —— 编辑对内容、来源或重建限度的评论。\textbf{[无法辨认]} —— 原文因载体退化而无法破译的段落。\textbf{[缺文]} —— 原始笔记本中缺失的书页或段落。\textbf{[原文如此]} —— 对原文错误或特质的逐字转录。

\vspace{0.4em}
\noindent\textbf{限度:} 本版不声称完整。原始日记约有847页；读者拿到的是经过编选的选本，按五个主题维度（MEM、DOC、LUC、CONT）加尾声组织。本三语版将每个碎片以三种语言呈现（葡萄牙语 → 英语 → 简体中文），作为已出版各版的镜像。

\vspace{1.2em}

\begin{flushright}
--- A Edição --- The Edition --- 编委会\par
OpenCode Books, Crateús\par
Julho de 2026 --- July 2026 --- 2026年7月
\end{flushright}

\endgroup
\clearpage
